% Part: sets-functions-relations
% Chapter: sets
% Section: proofs-about-sets

\documentclass[../../../include/open-logic-section]{subfiles}

\begin{document}

\olfileid[pt]{sfr}{set}{prf}
\olsection{Provas sobre Conjuntos}

\begin{editorial}
Esta seção foi substituída por \verb|content/method/proofs| e foi
removida deste capítulo por padrão.
\end{editorial}

\begin{explain}
Conjuntos e as notações que introduzimos até aqui nos fornecem
abreviações convenientes para especificar conjuntos e expressar relações
entre eles. Com frequência também será necessário provar afirmações sobre
essas relações. Se você não está familiarizado com provas matemáticas,
isso pode ser novo para você. Então percorreremos um exemplo simples. Vamos
provar que para quaisquer conjuntos $X$ e $Y$, é sempre o caso que $X \cap
(X \cup Y) = X$. Como se prova uma identidade entre conjuntos assim?
Lembre que conjuntos são determinados unicamente por seus !!{element}s, isto é,
conjuntos são idênticos sse têm os mesmos !!{element}s. Então neste
caso temos que provar que (a) todo !!{element} de $X \cap (X \cup
Y)$ é também !!a{element} de $X$ e, reciprocamente, que (b) todo
!!{element} de $X$ é também !!a{element} de $X \cap (X \cup Y)$. Em
outras palavras, mostramos que tanto (a) $X \cap (X \cup Y) \subseteq X$ quanto
(b) $X \subseteq X \cap (X \cup Y)$.

Uma prova de uma afirmação geral como ``todo !!{element}~$z$ de $X \cap (X
\cup Y)$ é também !!a{element} de $X$'' é provada primeiro assumindo
que um $z \in X \cap (X \cup Y)$ arbitrário é dado, e provando a partir
dessa suposição que $z \in X$. Você pode conhecer esse padrão como ``prova
condicional geral.'' Nesta prova também teremos que usar as
definições envolvidas na suposição e conclusão, p.ex., neste
caso de ``$\cap$'' e ``$\cup$.'' Assim o caso (a) seria argumentado da
seguinte forma:

\begin{quote}
(a) Primeiro queremos mostrar que $X \cap (X \cup Y) \subseteq X$, isto é,
pela definição de $\subseteq$, que se $z \in X \cap (X \cup Y)$ então
$z \in X$, para qualquer~$z$. Assuma então que $z \in X \cap (X \cup
Y)$. Como $z$ é !!a{element} da interseção de dois conjuntos sse
é !!a{element} de ambos os conjuntos, podemos concluir que $z \in X$ e
também $z \in X \cup Y$. Em particular, $z \in X$, que é o que
queríamos mostrar.
\end{quote}

Isso completa a primeira metade da prova. Note que no último
passo usamos o fato de que se uma conjunção ($z \in X$ e $z \in X
\cup Y$) segue de uma suposição, cada conjunto segue dessa mesma suposição. Você pode conhecer essa regra como ``eliminação da conjunção,'' ou $\land$Elim. Agora vamos provar (b):

\begin{quote}
(b) Agora provamos que $X \subseteq X \cap (X \cup Y)$, isto é, pela
definição de $\subseteq$, que se $z \in X$ então também $z \in X \cap
(X \cup Y)$, para qualquer~$z$. Assuma $z \in X$. Para mostrar que $z \in X
\cap (X \cup Y)$, temos que mostrar (pela definição de ``$\cap$'') que
(i) $z \in X$ e também (ii) $z \in X \cup Y$. Aqui (i) é apenas nossa
suposição, então não há nada mais a provar. Para (ii), lembre que $z$
é !!a{element} de uma união de conjuntos sse é !!a{element} de pelo menos
um desses conjuntos. Como $z \in X$, e $X \cup Y$ é a união de
$X$ e $Y$, este é o caso aqui. Então $z \in X \cup Y$. Mostramos
tanto (i) $z \in X$ quanto (ii) $z \in X \cup Y$, logo, pela definição de
``$\cap$,'' $z \in X \cap (X \cup Y)$.
\end{quote}

Isso foi um pouco prolixo, mas ilustra como raciocinamos sobre
conjuntos e suas relações. Geralmente não somos tão explícitos; em
particular, podemos não repetir todas as definições. Uma prova do nosso
resultado em um texto mais avançado seria muito mais comprimida. Poderia
parecer algo assim.
\end{explain}


\begin{prop}[Absorção]
Para todos os conjuntos $X$, $Y$,
\[
X \cap (X \cup Y) = X
\]
\end{prop}

\begin{proof}
(a) Suponha $z \in X \cap (X \cup Y)$. Então $z \in X$, logo $X \cap (X
  \cup Y) \subseteq X$.

(b) Agora suponha $z \in X$. Então também $z \in X \cup Y$, e portanto
  também $z \in X \cap (X \cup Y)$. Assim, $X \subseteq X \cap (X \cup
  Y)$.
\end{proof}

\begin{prob}
Prove em detalhe que $X \cup (X \cap Y) = X$. Depois dê uma prova
comprimida e mais curta. (Dica: para a direção $X \cup (X \cap Y) \subseteq X$
você precisará de prova por casos, aka $\lor$Elim.)
\end{prob}

\end{document}
